% mazzi: 01
\chapter{Cenni di storia dell'anatomia}

% ---------------------------------------------------------------------------
\section{Definizione di anatomia}

\textbf{Anatomia} (gr.\ \textit{anatomè} = ``dissezione'', da \textit{anà} ``attraverso'' +
\textit{tèmno} ``tagliare''): branca della biologia che studia \textbf{forma e struttura} degli
organismi.

\rubrica{Suddivisioni}
\begin{itemize}
  \item \textbf{normale} e \textbf{patologica};
  \item artistica, superficiale, di comparazione anatomica, esperienziale.
\end{itemize}
Anatomia e fisiologia studiano insieme la forma e il funzionamento del corpo.

% ---------------------------------------------------------------------------
\section{Breve storia dell'anatomia}

\subsection{L'Antico Egitto}
Ordine descrittivo geografico: dall'alto (capo) fino ai piedi. Conoscenze documentate da papiri:
\begin{itemize}
  \item \textbf{1850 a.C.} --- Papiro di Kahun: ginecologia;
  \item \textbf{1550 a.C.} --- Papiro di Ebers (anatomia e patologie), Papiro di Hearst (apparato
    urinario, sangue, denti);
  \item documenti su chirurgia delle fratture e medicamenti.
\end{itemize}
\textit{Libro dei Morti di Hunefer} (British Museum): descrive la cerimonia di apertura della
bocca $\rightarrow$ conoscenza pratica del corpo legata al rito funerario.

\subsection{La Grecia antica e il Medioevo}
\begin{itemize}
  \item primi esperimenti anatomici sistematici in Grecia;
  \item \textbf{dissezione dei cadaveri} già nel XIII sec.: prima come esame autoptico (morti in
    circostanze dubbie), poi a scopo didattico;
  \item \textbf{Bologna} = prima sede di rilievo: testo di chirurgia di Guglielmo da Saliceto
    (1270, osservazione diretta) $\rightarrow$ \textbf{\textit{Anathomia Mundini}} (1316, Mondino
    dei Liuzzi), il trattato più importante del Medioevo;
  \item \textbf{scuola anatomica di Padova} $\rightarrow$ \textbf{Andrea Vesalius},
    \textit{De humani corporis fabrica} (\textbf{1543}) = nascita dell'Anatomia Moderna
    (Rinascimento), basata sull'osservazione diretta su cadaveri.
\end{itemize}

\rubrica{Medicina nel Medioevo}
\begin{itemize}
  \item pratica e diretta;
  \item cure con stregonerie, incantesimi, terapie a base d'erbe;
  \item nascita dei primi ospedali, ospizi e monasteri.
\end{itemize}
\textit{Wound man}: figura anatomica usata come guida per i chirurghi (o barbieri-chirurghi) sulle
potenziali ferite, per rivelare cure e trattamenti disponibili (es.\ manoscritto fine XV sec.,
Bayerische Staatsbibliothek Cgm 597). Strumenti chirurgici del XV sec.: forbici, forbice-pinza,
uncino, forcella, sonda, bisturi, aghi, filo, sega.

\subsection{Mondino de' Liuzzi}
Tra i più grandi anatomisti italiani (attivo $\sim$1300). Usò cadaveri di criminali, in lotta col
rapido deterioramento dei corpi. Nel \textbf{1316} pubblicò il manuale \textit{Anothomia}.

\subsection{Il Rinascimento: Leonardo da Vinci}
Svolta fondamentale: l'occhio dell'artista vede linee, volumi e proporzioni.
\begin{quote}
\textit{``\ldots questa mia figurazione del corpo umano ti sarà dimostrata non altrementi che se
tu avessi l'omo naturale innanzi.''}\\
\hfill --- Leonardo da Vinci
\end{quote}
\begin{itemize}
  \item vera innovazione = \textbf{ricerca della funzione} degli organi $\rightarrow$ studio
    fisiologico applicato a quello anatomico;
  \item \textbf{vedute esplose}: scompongono i particolari nello spazio per comprenderli meglio
    $\rightarrow$ impulso all'\textbf{anatomo-fisiologia};
  \item dissezioni in prima persona (famosa quella del ``vecchio di cent'anni'', Ospedale di
    Santa Maria Nuova, Firenze);
  \item disegni anatomici (cranio, colonna, arti, organi interni, feto) tra $\sim$1480 e 1517.
\end{itemize}

\subsection{Guido da Vigevano}
Trattato \textit{Anothomia designata per figuras} (1345): uno dei primi esempi di anatomia
illustrata. Confrontato con Leonardo alla Biblioteca Ambrosiana di Milano (6 febbraio 2020): 18
figure (in 16 tavole) accostate ai disegni anatomici di Leonardo.

\subsection{L'epoca moderna: dal 1800 ad oggi}
\begin{itemize}
  \item \textbf{1800}: trafugamento di cadaveri, sezionamento di soggetti vivi, abusi $\rightarrow$
    inizio della storia dell'\textbf{Anatomia Patologica};
  \item si afferma lo studio della \textbf{Morfologia}: forma delle strutture e rapporti tra organi;
  \item oggi l'anatomia studia a più \textbf{livelli di osservazione}: tessuti, cellule, molecole.
\end{itemize}
Il corpo umano = sistema biologico di organi e apparati $\rightarrow$ tessuti $\rightarrow$ cellule
+ tessuto connettivo. La \textbf{Anatomia del Gray} (\textit{Anatomy of the Human Body}, Henry
Gray) è il riferimento per l'apprendimento della disciplina, base della pratica clinica.

% ---------------------------------------------------------------------------
\section{Anatomia microscopica: citologia e istologia}

\subsection{Anatomia microscopica}
\textbf{Anatomia microscopica}: studia le strutture non visibili senza ingrandimento. Si articola in:
\begin{itemize}
  \item \textbf{citologia}: analizza i singoli componenti cellulari;
  \item \textbf{istologia}: esamina i tessuti (insieme di cellule e prodotti cellulari che
    cooperano per una o più funzioni specifiche).
\end{itemize}
Tutte le cellule del corpo si suddividono in \textbf{4 tipi tessutali}:
\begin{enumerate}
  \item epiteli;
  \item tessuti connettivi;
  \item tessuto muscolare;
  \item tessuto nervoso.
\end{enumerate}

\subsection{Istologia}
Branca della biologia applicata alla medicina, con ruolo centrale anche in Anatomia Patologica.
\begin{itemize}
  \item \textbf{Albrecht von Haller} (1708--1777): primo a suggerire un'organizzazione strutturale
    microscopica degli esseri viventi;
  \item \textbf{istologia moderna}: nasce nel XIX sec.\ con la \textbf{teoria cellulare}.
\end{itemize}
Metodo: microscopio $\rightarrow$ i tessuti sono tagliati in strisce sottilissime, trattati per
prevenire la decomposizione e colorati per rendere visibili le strutture.

% ---------------------------------------------------------------------------
\section{Organizzazione del corpo umano}

Due livelli principali di studio:
\begin{itemize}
  \item \textbf{organizzazione macroscopica}: imaging (radiografia, TAC, RMN, ecografia, PET)
    $\rightarrow$ esplora l'interno del corpo senza dissezione;
  \item \textbf{organizzazione microscopica}: istologia e citologia.
\end{itemize}
Esempio contemporaneo: mostra \textit{Real Bodies} (Milano), cadaveri plastinati esposti per
illustrare l'organizzazione interna del corpo.

% ---------------------------------------------------------------------------
\section{L'anatomia e il futuro}

Disciplina in continua evoluzione: dall'osservazione macroscopica dei cadaveri all'imaging
digitale tridimensionale, all'integrazione con la biomeccanica, fino alla bioingegneria (protesi
neurali, arti artificiali). Anatomia e fisiologia restano le basi irrinunciabili della
comprensione del corpo umano e della pratica clinica.
